\section{Induktion\label{KapInduktion}}
\begin{frame} \ftx{Induction}
	\index{Induktion}
	\s{
		Induction describes the generation of an electrical voltage in a conductor due to a change in magnetic flux over time. The change in magnetic flux can be caused in various ways.
        Figure \ref{Induktion1} illustrates induction using a moving conductor in a magnetic field. As in the previous example (Figure \ref{Lorentz1}), a conductor is penetrated by the magnetic field of a permanent magnet, but with the difference that this time no current flows through the conductor. 
 		Instead, the ends of the conductor are connected to a voltage measuring device. If the conductor is now moved by an external force, the Lorentz force acts on the electrons in the conductor. The movement causes a charge separation in the conductor. This charge separation is called induction voltage. As long as the conductor or the inducing magnetic field is in motion, a voltage can be read on the measuring device.
		\fo{width=0.45\textwidth}{}{Induktion_neu}{
			\put(15, 35){\color{gray}$v$}
			\put(78, 72){\color{magnetfeld}$N$}
			\put(78, 33){\color{magnetfeld}$S$}
			\put(15, 23){\color{magnetfeld}$\vec{B},\,\varPhi $}
			\put(49, 39){\color{gray}$\ell$}
			\put(17, 73){\input{Tikz/src/messung_induzierte_spannung}
			}
		}{{\bf Experimental setup for measuring an induced voltage.} If a time-varying magnetic field is generated by the movement of the conductor, the Lorentz force causes charge separation in the conductor, inducing a measurable voltage. \label{Induktion1}}

        In simple terms, the induced voltage can be calculated using equation \ref{GlvereinfachtInduktionsgesetz}. The magnitude of the induced voltage $u_{\mathrm{i}}$ depends on the magnetic flux density $B$, the length of the conductor in the magnetic field $\ell$, the speed of movement or flux change $v$, and the number of conductors in the magnetic field $N$. Since this voltage varies over time, the symbol $u_{\mathrm{i}}$ is written in lower case.			\begin{eq}
				u_{\mathrm{i}} = B\cdot \ell\cdot v\cdot N \label{GlvereinfachtInduktionsgesetz}
			\end{eq}\\
	
		In general, the induced voltage $u_{\mathrm{i}}$ is the change in magnetic flux ${\d \varPhi}/{\d t}$ over time multiplied by the number of conductors N. This is reflected in the general law of induction (equation \ref{GlInduktionsgesetz}).
        The negative sign is due to Lenz's law, which states that cause and effect always behave in opposite ways.
	}
	
	\b{
		\begin{minipage}{0.4\textwidth}
			\fo{}{width=\textwidth}{Induktion_neu}{
				\put(15, 34){\color{gray}$v$}
				\put(78, 72){\color{magnetfeld}$N$}
				\put(78, 33){\color{magnetfeld}$S$}
				\put(11, 23){\color{magnetfeld}$\vec{B},\,\varPhi $}
				\put(49, 39){\color{gray}$\ell$}
				\put(17, 73){\begin{tikzpicture}[scale=0.73]
					\draw (0,0) -- ++ (1.3,-0.65) -- ++(0,0.7) coordinate (u1);
					\draw (3.6,-1.8) -- ++ (-1.3,0.65) -- ++(0,0.7) coordinate (u2);
					\draw (u2) to[open, o-o] (u1);
					\draw[->, color=voltage, thick] (2.15, -0.38) -- (1.45, 0);
					\draw [color=voltage] (2, -0.2) node[above] {$u_{\mathrm{i}}$};
				\end{tikzpicture}
				}
			}{}
		\end{minipage}%
		\pause%
		\hfill%
		\begin{minipage}{0.57\textwidth}
			\begin{eq}
				u_{\mathrm{i}} = B\cdot \ell\cdot v\cdot N
			\end{eq}\\
			Die Höhe der induzierten Spannung hängt ab von:
			\begin{itemize}
				\item der magnetischen Flussdichte $B$
				\item der Länge des Leiters im Magnetfeld $\ell$
				\item der Geschwindigkeit der Bewegung oder \\Flussänderung $v$
				\item der Anzahl der Leiter im Magnetfeld $N$
			\end{itemize}
		\end{minipage}
	}
	% \f{trim=0 0cm 0 0,clip,width=0.5\textwidth}{trim=0 0cm 0 0,clip,width=0.45\textwidth}{Induktion1}{Messung der induzierten Spannung \label{Induktion1}{\tiny(Quelle: \cite[Kapitel 5.6.1]{Tkotz})}}

\end{frame}
\begin{frame} \ftx{Induktion}

	\b{Allgemeines Induktionsgesetz:}
	\begin{eq}
		u_{\mathrm{i}} = -N\cdot \frac{\d \varPhi}{\d t}\label{GlInduktionsgesetz}
	\end{eq} 
	\pause
\s{\begin{Key point}{Induction}
    Induction is the electrical voltage $u_{\mathrm{i}}$ generated by a change in magnetic flux ${\d \varPhi}/{\d t}$.
\end{Key point}}
	\b{Lenzsche Regel\index{Lenzsche Regel}:
	\par	
	\center{Ursache und Wirkung sind stets entgegengesetzt.}
	}
\end{frame}